\section{Evaluation y Tooling: primero prevenir la Contamination, después hablar de Leaderboards}

Que termine el preentrenamiento no significa que los resultados sean confiables. Los benchmarks de código tienden especialmente a aparecer en GitHub, en tutoriales y en instruction data; además, los leaderboards públicos incentivan la optimización repetida. Por ello, el curso propone usar múltiples benchmarks, decontamination, membership test y time-split evaluation.

\subsection{Múltiples benchmarks: una sola puntuación no describe un code model}

La primera tabla compara al mismo tiempo code completion, code reasoning, data science, math y repository tasks. Un modelo de mayor size suele ser más fuerte, pero el orden de los modelos puede variar entre tasks. La evaluación debe ir ligada a model version, prompt, sampling, tests y execution environment.

\lecturefigure{slide-64.jpg}{Comparación en múltiples benchmarks de StarCoder/CodeLlama/DeepSeek y otros modelos}{deck oficial de Loubna Ben Allal, página 64}

\lecturefigure{slide-65.jpg}{Resultados desglosados por tarea de StarCoder2 3B/7B/15B}{deck oficial de Loubna Ben Allal, página 65}

Si para cada problema se muestrean $n$ solutions y $c$ de ellas pasan los tests, el estimador insesgado habitual de pass@$k$ es
\[
\operatorname{pass@}k=1-\frac{\binom{n-c}{k}}{\binom{n}{k}}.
\]
pass@1 se aproxima más a una sola completion; un $k$ mayor mide que al menos uno de varios intentos tenga éxito, por lo que los valores no pueden compararse directamente.

\begin{warningbox}{Los unit tests solo cubren el comportamiento que se escribió}
Pasar los tests no garantiza efficiency, security, style, ausencia de undefined behavior, dependency correctness ni el manejo de hidden edge cases. Si los tests son demasiado débiles, el modelo puede generar código que parece correcto pero no es confiable.
\end{warningbox}

\teachervoice{00:44:20--00:46:20, el ponente recomienda usar en el release tantos benchmarks como sea posible: alguno puede tener contamination, y las múltiples tareas también revelan la distribución del comportamiento del modelo, en lugar de perseguir solo un headline individual.}

\subsection{Autocomplete y membership test: el lado del despliegue también debe hacer attribution}

BigCode también publicó una VS Code extension y un membership test. Este último intenta determinar si un generated snippet se parece a los training data y advierte al usuario sobre su posible origen. Es un provenance aid, no una memorization proof concluyente.

\lecturefigure{slide-66.jpg}{Code tooling: autocomplete extension y training-data membership test}{deck oficial de Loubna Ben Allal, página 66}

\begin{knowledgebox}{Límites del membership test}
El exact match tiene alta precisión pero pasa por alto las reescrituras; el fuzzy match marca common idioms como falsos positivos; un training corpus incompleto o de otra versión produce false negatives. La salida debe ser una similitud y unas fuentes candidatas, no redactarse como una conclusión legal.
\end{knowledgebox}

\teachervoice{00:45:00--00:47:00, el ponente ubica el membership test dentro de los esfuerzos de code attribution: cuando el código generado puede provenir de ejemplos de entrenamiento, se debe dar al autor o al usuario una advertencia que pueda verificar.}

\subsection{Customize code model: cómo llevar los principios del preentrenamiento a una sola GPU}

El Personal Copilot workflow sigue siendo collect--filter--dedup--format--train--evaluate; lo único que cambia es que el full pretraining se sustituye por repository-specific continued pretraining o fine-tuning. LoRA expresa la actualización de los pesos como una matriz de bajo rango
\[
W'=W+BA,\qquad A\in\mathbb{R}^{r\times d},\ B\in\mathbb{R}^{d'\times r},\ r\ll d,
\]
y solo entrena $A,B$, lo que reduce notablemente la memory de los optimizer states y de los gradientes.

\lecturefigure{slide-67.jpg}{Customize Code Models: entrenar un personal Copilot sobre un codebase privado}{deck oficial de Loubna Ben Allal, página 67}

\begin{lstlisting}[caption={Ruta ejecutable para un code assistant con presupuesto reducido}]
select_small_strong_base_with_code_and_long_context()
collect_only_authorized_repository_data()
deduplicate_and_split_by_repository_or_time()
format_fim_and_repository_context_examples()
fine_tune_with_peft_or_lora_on_one_gpu()
evaluate_execution_latency_leakage_and_regressions()
\end{lstlisting}

\teachervoice{00:57:20--00:59:10, en la sesión de Q\&A el ponente da recomendaciones para una sola GPU: elegir una base fuerte, curated data, un runtime quantized/on-device y PEFT, en lugar de intentar preentrenar desde cero.}

\subsection{Leaderboards: la interfaz es cómoda, pero el protocol importa más}

Los leaderboards públicos agregan model scores y ayudan a identificar candidatos; pero detrás de cada row hay model date, prompt, temperature, test harness, contamination policy y submission rules. La captura de pantalla es un punto de entrada, no un registro de experimento reproducible.

\lecturefigure{slide-68.jpg}{Code leaderboards: punto de entrada público para comparar modelos}{deck oficial de Loubna Ben Allal, página 68}

\begin{knowledgebox}{Orden de uso de un leaderboard}
Primero revisa la benchmark version y el cutoff; después, la release/training date del modelo; confirma base/instruct, pass@1/pass@k, execution sandbox y prompt; solo al final compara las puntuaciones. Si el protocol es distinto, la clasificación carece de interpretabilidad.
\end{knowledgebox}

\subsection{HumanEval overfitting: el instruction tuning también contamina}

Aunque el pretraining elimine HumanEval, los instruction data pueden contener ejercicios parecidos del tipo «implementar una función y pasar los tests». El modelo puede aprender el benchmark style o haber visto problemas semánticamente equivalentes, lo que infla el public score respecto de problemas realmente nuevos.

\lecturefigure{slide-69.jpg}{Potential overfitting in HumanEval: brecha entre el benchmark público y la generalización real}{deck oficial de Loubna Ben Allal, página 69}

Si la similarity entre el ejemplo de entrenamiento $x$ y el benchmark $b$ cumple $s(x,b)>\tau$, puede marcarse un posible overlap; sin embargo, un lexical detector no necesariamente capta la semantic equivalence. La decontamination debe combinar exact/fuzzy/code-structure/time metadata y reportar los umbrales.

\begin{warningbox}{Los instruction data son la segunda vía de entrada de la contamination}
No basta con limpiar el pretraining corpus. Los SFT, los RLHF prompts, los synthetic exercises, las benchmark explanations y las GitHub solutions pueden filtrar el test pattern; cada stage requiere un provenance scan.
\end{warningbox}

\subsection{LiveCodeBench: aislamiento temporal mediante la release date}

LiveCodeBench recopila periódicamente nuevos contest problems y evalúa solo con los problemas posteriores al release del modelo, lo que reduce la probabilidad de que este los haya visto durante el entrenamiento. También compara el full-history score con el post-cutoff score; si la diferencia es grande, puede indicar overfit a benchmarks antiguos.

\lecturefigure{slide-70.jpg}{LiveCodeBench: contamination-resistant evaluation con post-release problems}{deck oficial de Loubna Ben Allal, página 70}

La time-split puede definirse como
\[
\mathcal{B}_{\text{valid}}(m)=\{b:\ t_{\text{problem}}(b)>t_{\text{release}}(m)\}.
\]
Reduce la fuga temporal directa, pero el modelo puede basarse en un checkpoint no publicado, en problemas que aparecieron antes en la web o en templates similares, por lo que se requiere cautela.

\teachervoice{00:49:00--00:52:40, el ponente presenta LiveCodeBench como una solución importante: usar problemas que aparecieron después de la publicación del modelo y observar si algunos modelos mantienen su desempeño en el uncontaminated slice.}

\subsection{Resumen de la sección}

Una code evaluation confiable requiere múltiples tareas, tests ejecutables, un pass@k explícito, decontamination de datos y temporal, herramientas de membership/attribution y un repeated protocol. El leaderboard sirve para orientarse; la time-split y el source audit son la defensa real contra la contamination.
